\documentclass[11pt,a4paper]{article}

\usepackage[utf8]{inputenc}
\usepackage[ngerman]{babel}
\usepackage[T1]{fontenc}
\usepackage{geometry}
\geometry{a4paper, left=25mm, right=25mm, top=25mm, bottom=25mm}
\usepackage{graphicx}
\usepackage{xcolor}
\usepackage{tabularx}
\usepackage{enumitem}
\usepackage{hyperref}
\usepackage{listings}
\usepackage{booktabs}
\usepackage{tcolorbox}
\usepackage{fancyhdr}

% Farben definieren
\definecolor{siemensblue}{RGB}{0, 153, 153}
\definecolor{lightgray}{RGB}{240, 240, 240}
\definecolor{warningred}{RGB}{204, 0, 0}

% Kopf- und Fußzeile
\pagestyle{fancy}
\fancyhf{}
\fancyhead[L]{\textbf{Laborübung: Automatisierungstechnik}}
\fancyhead[R]{OPC UA \& S7-1500}
\fancyfoot[C]{\thepage}

% Hyperref Setup
\hypersetup{
    colorlinks=true,
    linkcolor=siemensblue,
    filecolor=magenta,      
    urlcolor=siemensblue,
    pdftitle={Laborübung Rolltor-Simulation},
}

\begin{document}

\begin{center}
    \vspace*{1cm}
    \Huge \textbf{Laboraufgabe: Automatisierung eines Rolltors via OPC UA} \\[0.5cm]
    \Large Praktikum Steuerungstechnik mit Siemens S7-1500 \\[1.5cm]
\end{center}

\section*{Zielsetzung der Laborübung}
In dieser Laborübung programmieren Sie die Steuerung für eine Rolltor-Simulation. Die physikalische Simulation des Tors läuft als eigenständiges Python-Programm und kommuniziert über das standardisierte Protokoll \textbf{OPC UA} mit einer simulierten Speicherprogrammierbaren Steuerung (Siemens S7-1500 über PLCSIM Advanced). 

Lernziele dieser Übung sind:
\begin{itemize}[noitemsep]
    \item Einrichtung eines OPC UA Servers auf einer S7-1500.
    \item Erstellung von Datenaustausch-Schnittstellen (Datenbausteinen).
    \item Programmierung grundlegender Verriegelungen und Selbsthaltungen in KOP und FUP.
    \item Fehlerbehandlung und Ablaufsteuerungen in SCL.
\end{itemize}

\vspace{0.5cm}
\begin{tcolorbox}[colback=warningred!10, colframe=warningred, title=\textbf{Wichtiger Hinweis zur Simulation}]
Die Rolltor-Simulation simuliert reale physikalische Grenzen. Wird ein Motor nicht rechtzeitig gestoppt und überfährt das Tor die Endlagen ($> 102\,\%$ oder $< -2\,\%$), geht das System in einen kritischen Fehlerzustand (Overtravel Fault). Gleiches gilt, wenn beide Motoren (Auf/Ab) gleichzeitig angesteuert werden.
\end{tcolorbox}

\newpage

\section{Phase 1: Infrastruktur und Setup}

Bevor Sie mit der eigentlichen Programmierung beginnen können, muss die Kommunikationsschnittstelle zwischen dem TIA Portal (SPS) und der Simulation (Python) eingerichtet werden.

\subsection{Aufgabe 1.1: Datenbaustein (DB) anlegen}
Legen Sie im TIA Portal einen neuen globalen Datenbaustein an (z.\,B. \texttt{DB\_Rolltor\_OPC}). Dieser dient als Schnittstelle. Legen Sie darin exakt die folgenden Variablen an:

\begin{table}[h]
    \centering
    \renewcommand{\arraystretch}{1.2}
    \begin{tabularx}{\textwidth}{@{} l l X c @{}}
        \toprule
        \textbf{Variablenname} & \textbf{Datentyp} & \textbf{Beschreibung / Richtung} & \textbf{Erreichbar (OPC)} \\
        \midrule
        \texttt{Position} & \texttt{Real} & Aktuelle Position des Tors (0.0 bis 100.0) [Eingang] & Ja \\
        \texttt{Error} & \texttt{Bool} & Störsignal der Simulation (True = Fehler) [Eingang] & Ja \\
        \texttt{Upper\_Limit} & \texttt{Bool} & Oberer Endschalter ($100\,\%$) erreicht [Eingang] & Ja \\
        \texttt{Lower\_Limit} & \texttt{Bool} & Unterer Endschalter ($0\,\%$) erreicht [Eingang] & Ja \\
        \texttt{Motor\_Up} & \texttt{Bool} & Steuerbefehl: Motor AUF [Ausgang] & Ja \\
        \texttt{Motor\_Down} & \texttt{Bool} & Steuerbefehl: Motor AB [Ausgang] & Ja \\
        \bottomrule
    \end{tabularx}
\end{table}

\subsection{Aufgabe 1.2: OPC UA Server aktivieren}
\begin{enumerate}[noitemsep]
    \item Öffnen Sie die Gerätekonfiguration der S7-1500.
    \item Navigieren Sie zu \textit{OPC UA} $\rightarrow$ \textit{Server} und aktivieren Sie den OPC UA Server.
    \item Notieren Sie sich die Server-URL (z.\,B. \texttt{opc.tcp://192.168.0.1:4840}).
    \item Stellen Sie sicher, dass in den Eigenschaften des erstellten Datenbausteins die Variablen für OPC UA freigegeben sind (Spalte: \textit{Erreichbar aus HMI/OPC UA}).
    \item Übersetzen Sie das Projekt und laden Sie es in die simulierte Steuerung (PLCSIM Advanced).
\end{enumerate}

\subsection{Aufgabe 1.3: Verbindungstest}
\begin{enumerate}[noitemsep]
    \item Verbinden Sie sich zunächst mit einem Diagnose-Tool (z.\,B. \textit{UaExpert}) auf die SPS und prüfen Sie, ob Sie die Variablen im Adressraum sehen und lesen/schreiben können.
    \item Tragen Sie die korrekten Node-IDs (z.\,B. \texttt{ns=3;s="DB\_Rolltor\_OPC"."Motor\_Up"}) und die Server-URL in die \texttt{config.json} der Python-Simulation ein.
    \item Starten Sie die Simulation im \texttt{client}-Modus und prüfen Sie, ob sich die Werte bei Betätigung der GUI-Buttons im TIA Portal ändern.
\end{enumerate}

\newpage

\section{Phase 2: Programmieraufgaben (IEC 61131-3)}

Lösen Sie die folgenden Aufgaben in den angegebenen Programmiersprachen. Kapseln Sie jede Aufgabe sinnvoll in eigenen Funktionen (FC) oder Funktionsbausteinen (FB) und rufen Sie diese im \texttt{OB1} (Main) auf.

\subsection{Aufgabe 2.1: Totmannsteuerung (Tippbetrieb)}
\textbf{Sprache:} Kontaktplan (KOP) \\
\textbf{Szenario:} Das Tor soll manuell gesteuert werden können, solange Taster gedrückt gehalten werden.
\begin{itemize}
    \item Erstellen Sie zwei Merker/Eingänge in der SPS für virtuelle Taster: \texttt{Taster\_Hand\_Auf} und \texttt{Taster\_Hand\_Ab}.
    \item Das Tor soll nur fahren, solange der jeweilige Taster ein \texttt{True}-Signal liefert.
    \item \textbf{Verriegelung:} Es muss hardwaretechnisch ausgeschlossen sein, dass \texttt{Motor\_Up} und \texttt{Motor\_Down} gleichzeitig aktiv sind.
    \item \textbf{Endlagen-Stopp:} Wird der Endschalter (\texttt{Upper\_Limit} bzw. \texttt{Lower\_Limit}) erreicht, muss der jeweilige Motor zwingend abschalten, auch wenn der Taster noch gedrückt ist.
\end{itemize}

\subsection{Aufgabe 2.2: Automatikbetrieb mit Selbsthaltung}
\textbf{Sprache:} Funktionsplan (FUP) \\
\textbf{Szenario:} Das Tor soll nach einem kurzen Tastendruck automatisch bis zur Endlage fahren.
\begin{itemize}
    \item Nutzen Sie SR-Flipflops (Setzen/Rücksetzen), um den Motor-Zustand zu speichern.
    \item Ein kurzer Impuls auf \texttt{Taster\_Auto\_Auf} soll das Flipflop für \texttt{Motor\_Up} setzen.
    \item \textbf{Wichtig:} Das Rücksetzen (Reset) des Flipflops muss durch den Endschalter (\texttt{Upper\_Limit}) erfolgen. 
    \item Eine gegenseitige Verriegelung ist auch hier zwingend erforderlich!
    \item \textit{Achtung:} Wenn das Rücksetzen fehlschlägt, fährt die Simulation das Tor über $100\,\%$ hinaus in den mechanischen Crash (Overtravel Fault).
\end{itemize}

\subsection{Aufgabe 2.3: Störungsbehandlung und Diagnose}
\textbf{Sprache:} Structured Control Language (SCL) \\
\textbf{Szenario:} Die Anlage soll sicher auf Fehler reagieren und sich ordnungsgemäß quittieren lassen.
\begin{itemize}
    \item Wenn das Signal \texttt{Error} von der Simulation auf \texttt{True} geht, müssen alle Fahrbefehle (\texttt{Motor\_Up}, \texttt{Motor\_Down}) \textbf{sofort} überschrieben und auf \texttt{False} gesetzt werden.
    \item Dieser Störungs-Stopp hat die höchste Priorität im Programm.
    \item Implementieren Sie einen virtuellen Taster \texttt{Taster\_Reset}.
    \item Eine erneute Fahrt darf erst wieder freigegeben werden, wenn \texttt{Error = False} ist und der Fehler über \texttt{Taster\_Reset} (positive Flanke) quittiert wurde.
\end{itemize}

\subsection{Aufgabe 2.4: Automatische Zufahrt (Ablaufsteuerung)}
\textbf{Sprache:} SCL oder S7-GRAPH \\
\textbf{Szenario:} Das Tor soll im Automatikbetrieb nach einer Wartezeit selbstständig schließen.
\begin{itemize}
    \item Das Tor wird geöffnet. Sobald der obere Endschalter (\texttt{Upper\_Limit}) erreicht wird, soll eine Timer-Funktion (Einschaltverzögerung TON) gestartet werden.
    \item Die Wartezeit beträgt $10$ Sekunden.
    \item Nach Ablauf der 10 Sekunden schließt sich das Tor automatisch (Setzen von \texttt{Motor\_Down}).
    \item \textbf{Sicherheitsabbruch:} Wird während der Schließfahrt (oder während der Wartezeit) der Taster \texttt{Taster\_Auto\_Auf} gedrückt, muss die Zufahrt abgebrochen und das Tor sofort wieder geöffnet werden.
\end{itemize}

\vspace{1cm}
\begin{tcolorbox}[colback=siemensblue!10, colframe=siemensblue, title=\textbf{Abgabe}]
Senden Sie am Ende des Praktikums das archivierte TIA-Projekt (.zap16/.zap17) sowie ein kurzes PDF-Dokument mit Screenshots Ihrer Code-Bausteine und einer kurzen Beschreibung Ihrer Verriegelungs-Logik an den Betreuer.
\end{tcolorbox}

\end{document}
